\section{Stationary Reflected Comparisons}\label{s:coupling}

We prove the local comparison estimate used to pass from the constrained
microscopic dynamics to a continuum equation. The input consists of the
sparse-box estimates \ref{f:sparse} and \ref{f:conditioned-sparse}, the
extension over their holes in \ref{f:hole}, and the average bounds in
\ref{f:averages}. The principal new
point is a lower bound on dissipation obtained by forcing a contact inside
a fixed box. Its constants remain uniform when a constraint becomes
redundant, and after adapted exact-value observations.

\subsection{Experiments and the comparison statement}

We use the unit triangular lattice, with macroscopic scale $n$. Write
$Q_r=x_0+[-rn,rn]^2$, using any fixed pair of lattice coordinates, and
$\widehat Q_r=n^{-1}Q_r$ for its scaled version. For a vertex function,
$\int_{n,Q}F=n^{-2}\sum_{x\in Q}F(x)$; for an edge function the same
notation sums over unoriented lattice edges with midpoint in $Q$.
Omitting $Q$ means the currently specified localization box.
Constants may depend on fixed box ratios but not on $n$ or on the
diameter of the ambient finite domain. After choosing an orientation
$e=(x,y)$, a lattice gradient is the unscaled difference
$\nabla_e f=f(y)-f(x)$.
The notation $\Prob_{\rm vol}$ denotes the annealed law obtained by
first sampling the experiment and then an independent uniform vertex
in the indicated box, or an independent uniform edge for an edge
quantity. In a vertex-sampled gradient one fixes a lattice step
$\mathbf e$ and uses the edge $(X_n,X_n+\mathbf e)$.

An experiment starts with two finite-volume Lipschitz fields. Each has
bounded Dirichlet data and finitely many prescribed sign strings of the
type covered by \ref{f:sparse} and \ref{f:averages}. On the common free
graph in $Q_2$ their individual sign walls, when present, agree. We allow
the following additional observations. Starting from independent samples,
an algorithm queries both heights at a common sequence of vertices, with
each query and its stopping decision determined by the previous answers
and independent auxiliary randomness. Conditional on the complete
transcript, the queried heights are fixed. The two remaining conditional
laws are then put in a stationary synchronous joining. All estimates for
such an experiment are averaged over the original transcript as well as
the joining. The phrase ``uniform in the observations'' below has this
annealed meaning; it does not mean uniform over every possible exact-height
transcript.

We also allow jointly sampled boundary data when, conditional on those
data, the two interiors have their specified Gibbs kernels. Replacing
their conditional product by a stationary synchronous joining must
preserve each annealed marginal. The sparse-contact and average bounds
are required for those annealed marginals, with the same constants as
in the theorem below. This includes coupled random pins; it does not
assert a quenched estimate for an arbitrary realized pin vector.
The fixed set and the component on which a kernel is sampled may be
part of these recorded data. For example, a zero-pinned kernel on a
recorded component has the required bounds uniformly in that component,
provided its tested averages have the stipulated anchor. Its mixture
therefore has the same bounds. The other marginal must still satisfy
its own annealed estimates; selecting a component supplies no such
estimate by itself.

Let $P$ be the common fixed set in $Q_2$ and let $S$ be the set of common
individual sign walls. At time $t$ write $z_t=h_t-\widetilde h_t$.
For constants $a\le b$, if $a\le z|_P\le b$, set
\begin{equation}\label{c:tails}
 u^+_t=(z_t-b)_+,\qquad u^-_t=(a-z_t)_+.
\end{equation}
These functions have zero fixed input on $P$. When $S$ is nonempty we
take $a=b=0$ and use the two tails of $z$; square penalties at $S$ will
anchor their sum. In a region with no observations, we can instead use
$z$ modulo a spatial constant. Boundaries of $Q_2$ are artificial:
conditions outside that box need not agree.

The local average bounds required below are precisely those supplied by
\ref{f:averages}: averages of both marginals are bounded in every fixed
moment on the scale of $Q_2$ and its fixed-fraction subboxes, or differences
of equal-mass averages are so bounded in the free, centered case. The
constants for subbox side $\ell n$ may grow polynomially in $\ell^{-1}$.
For an adapted experiment these bounds concern the original marginal
laws; the conditioning argument below transfers them to the joined fields.
In a one-sided experiment we only require $z|_P\ge a$, and use $u^-$. An
observed boundary is \emph{controlled for the lower tail} if its fixed
values satisfy $z\ge a$; for the upper tail the condition is $z\le b$.
When extending a tail by zero behind such a boundary, we mean the side
removed from the common free graph, whose edges to the retained side
all pass through that observed boundary. An \emph{uncontrolled seed}
is a prescribed initial queried set on which no such inequality is
assumed. Any uncontrolled seed or other uncontrolled boundary must
remain a fixed positive scaled distance from the box on which an
estimate is asserted.

For a fixed step length $T$ and contact tolerance $\delta>0$, a
\emph{transmission edge} is a free--free lattice edge whose constraint
has slack at most $\delta$ at some time in $[0,T]$ in either replica.
The components of their union on the free vertices, including
singletons, are the transmission clusters. Write $C_x$ for the
cluster containing $x$ and $H_x=|C_x|$; also write $H_C=|C|$.
Contacts with fixed vertices or individual walls supply boundary
input to a cluster and do not connect clusters through a fixed
vertex. For a zero-drift cluster with zero tail input, its square
dissipation is
\[
 D_C=\sum_{x\in C}u_0(x)^2-\sum_{x\in C}u_T(x)^2.
\]
For two tails their dissipations are added. With drift, $D_C$ means
the sum of the nonnegative square losses at projection substeps,
with the source contribution kept separately. Lemma~\ref{c:cluster}
proves nonnegativity and the projection interpretation. Dissipation
sums concern clusters insulated from uncontrolled data; the
remaining clusters contribute the localization errors estimated
below. The theorem chooses $T$ and $\delta$ through its fixed
testing parameters.

\begin{theorem}[Local reflected comparison]\label{c:comparison}
Consider the experiments just specified, with $|a|+|b|$ bounded. Assume
the conclusions of \ref{f:conditioned-sparse}, \ref{f:hole}, and \ref{f:averages}
for their original marginal laws. Individual reflecting constraints are
ordinary one-sided sign walls, not pairs $0\le h_x\le\eps$ with
$\eps$ allowed to tend to zero. Fix $0<r<R<2$, $1<p<2$, and a desired
defect parameter $\rho>0$. There are fixed testing parameters, a step
length $T_\rho>0$, good edges $G_n$ and good wall sites $S_n^{\rm g}$,
and positive weights $w_e,w_x\le1$ with the following properties.

For either tail in \eqref{c:tails}, or for their signed difference, let
\begin{equation}\label{c:energy}
 \mathcal E_n(Q)=
 \sum_{e\in G_n,\ e\subset Q}w_e|\nabla_e u_0|^2
 +\sum_{x\in S_n^{\rm g}\cap Q}w_x|u_0(x)|^2.
\end{equation}
In the two-tail statement the energies are added. In a free comparison
the energy is that of $z$ and there is no wall sum. Then
\begin{equation}\label{c:energy-bound}
 \E\mathcal E_n(Q_r)\le C,\qquad
 \E\sum_{C:\,C\cap Q_r\ne\varnothing}D_C\le C.
\end{equation}
Here $C$ and $D_C$ have the preceding step-cluster meaning. The constants in
\eqref{c:energy-bound} may depend on $\rho,p,r,R$ and the average bounds.
Defective boxes satisfy the prescribed sparse-animal bound. For every
fixed $q<\infty$ and $A<\infty$, the normalized spatial sums of
$w^{-q}$, cluster-size moments, and hole-diameter moments are bounded by
deterministic constants outside an event of probability $O(n^{-A})$.
The relevant finite collection of moment orders can be chosen in advance.

Let $U_n$ be the extension over holes supplied by \ref{f:hole}, with a
fixed choice of the testing parameters. It can be set to zero on the
exceptional event. In scaled coordinates,
\begin{equation}\label{c:compactness}
 \E\|U_n\|_{W^{1,p}(n^{-1}Q_r)}^2\le C,\qquad
 \E\int_{n,Q_r}|U_n-u_0|\longrightarrow0.
\end{equation}
In a free comparison subtract a local cutoff-weighted average from both
$U_n$ and $z_0$. Its gradient and dissipation estimates are unchanged.
All these centered estimates are uniform in the size of the ambient
domain. The same bounds hold for a rough one-sided adapted experiment
with only $z|_P\ge a$, for $u^-=(a-z)_+$, and with zero extension behind
a controlled observed boundary. They hold uniformly after replacing $z$
by $z/K$, $K\ge1$, with bounded rescaled thresholds.

\emph{Threshold-uniform specialization.}
In the one-sided case, for each threshold $a$ satisfying the required
condition $z|_P\ge a$, the weights may be chosen independently of $a$.
With the original marginal moment bounds fixed,
\begin{equation}\label{c:threshold-scaling}
 \E\mathcal E_n\bigl((a-z)_+;Q_r\bigr)
 +\E\|U_n^{(a)}\|_{W^{1,p}(n^{-1}Q_r)}^2
 \le C_0(1+a^2).
\end{equation}
Thus, when $a=K\ge1$ and the observed values satisfy $z|_P\ge K$,
the energy and squared Sobolev norm of $(1-z/K)_+$ are bounded by
$2C_0$, uniformly in $K$.

For an ordinary simply connected polygon domain with identical fixed
boundary data, the exterior-preserving construction in
Lemma~\ref{c:exterior-zero} gives a bounded $W^{1,p}(\R^2)$ extension
which is exactly zero outside the lattice polygon. This construction
uses connectivity of its zero exterior; it is not asserted for an
arbitrary disconnected internal observed set.

If $Q_R$ has no observations or individual walls, and $X_n$ is sampled
from normalized ambient volume in $Q_r$, then
$n\nabla_{(X_n,X_n+\mathbf e)}z_0$ for each fixed lattice step $\mathbf e$,
$n(z_0(X_n+i)-z_0(X_n))$ for each fixed lattice offset $i$, and
$n\sup_{t\le T_0}|z_t(X_n)-z_0(X_n)|$ are tight. Thus this
untruncated conclusion has a fixed positive macroscopic free buffer.
It also holds up to controlled observations when $a=b$, because the
two tails then control $z-a$. For general $a<b$, corresponding
tightness statements hold only for the controlled tails up to
observations. Here $T_0$ is one fixed sufficiently short step.
Sampling is not conditioned on an arbitrarily small random free set.

Finally these statements allow a drift difference bounded by $c n^{-2}$,
with bounded support on the macroscopic scale, including a sufficiently
small smooth exponential tilt of either marginal. Constants may depend
on a fixed bound for $c$; the smallness needed for the change of marginal
law is that in \ref{f:averages}.
\end{theorem}

The wall sum in \eqref{c:energy} is retained for the subsequent trace
argument. It controls actual values at common walls, not a trace inferred
from interior convergence. The $L^1$ conclusion in \eqref{c:compactness}
will also be proved for a fixed hole construction; its holes need not
have density tending to zero.

\subsection{Conditioning and sparse exceptional events}

\begin{lemma}[Adapted freezing]\label{c:freezing}
Conditional on an adapted exact-value transcript, the unqueried fields
have the product of their Gibbs laws with those values pinned. Replacing
this product by a stationary synchronous coupling preserves each
annealed marginal and makes the coupled process stationary.
The same marginal-preservation statement holds for jointly sampled
boundary data with the conditional Gibbs kernels specified above.
\end{lemma}
\begin{proof}
For a fixed finite query history, every decision indicator is a function
of recorded values and auxiliary randomness. Disintegrate the original
product density over these values. The indicators introduce no factor
depending on an unqueried coordinate, so the remaining density is the
product of the two conditional Gibbs densities. Induction on the query
count, followed by partitioning over the stopping count, proves the
claim. Any coupling of these two kernels has the same individual
conditional marginals. Integration over the transcript recovers the
original marginals. A conditional invariant joining, obtained by time
averaging the synchronous semigroup on its compact state space, is
stationary. One may first take this time average jointly with the
transcript; disintegration of a weak limit supplies a measurable choice
of conditional invariant joinings.
For the boundary-data version, disintegrate directly over the joint
boundary vector and any recorded fixed set or component. Integrating
either marginal of a coupling of the two
conditional Gibbs kernels restores its original annealed law. The
subsequent sparse-event bounds use only these individual marginals;
independence of the two boundary vectors is unnecessary.
\end{proof}

Good-box marks must be recomputed on the initial pair of this joining.
Marks depending on unrevealed coordinates of the original independent
pair are not preserved. Their annealed estimates are preserved because
each rejoined marginal has its original law. If either marginal has a
$\rho^k$ bound on a separated $k$-set, the union of the two defective
sets has bound at most $2^k\rho^{k/2}$; this uses no independence of the
rejoined replicas.

\begin{lemma}[Conditional failures and moments]\label{c:conditional}
Let $A_i$ be future events and $\mathcal I$ the initial sigma-field.
If every subset $J$ of a separated set satisfies
$\Prob(\bigcap_{i\in J}A_i)\le\rho^{|J|}$, then, for $0<q<1$,
\[
 \Prob\left(\bigcap_{i\in I}
 \{\Prob(A_i\mid\mathcal I)>q\}\right)
 \le \frac{2}{q}\bigl(2\rho^{q/2}\bigr)^{|I|}.
\]
Also, if $0\le X_n\le n^B$ and
$\Prob(X_n>M)\le n^{-A}$, then
\[
 \Prob\{\E[X_n\mid\mathcal I]>M+1\}\le n^{B-A}.
\]
\end{lemma}
\begin{proof}
On the first event let $N=\sum_{i\in I}\mathbf1_{A_i}$ and $k=|I|$.
Then $\E[N\mid\mathcal I]>qk$ and $N\le k$, so
$\Prob(N\ge qk/2\mid\mathcal I)\ge q/(2-q)\ge q/2$.
The unconditional probability is at most
$2^k\rho^{\lceil qk/2\rceil}$ by choosing a witnessing subset.
For the second assertion use conditional Markov on
$\E[X_n\mathbf1_{X_n>M}\mid\mathcal I]$.
\end{proof}

In particular, for $q\ge1$, Jensen's inequality gives
\begin{equation}\label{c:conditional-moments}
 n^{-2}\sum_x\bigl(\E[H_x^k\mid\mathcal I]\bigr)^q
 \le\E\left[n^{-2}\sum_xH_x^{kq}\mid\mathcal I\right].
\end{equation}
Lower positive moment orders follow from any larger moment order.
The deterministic spatial moment bounds from \ref{f:hole}, with an
arbitrarily fast polynomial exceptional probability, pass to the
conditional moments in \eqref{c:conditional-moments}. First truncate
local clusters at a power of $n$ and discard the tail using the
arbitrarily high cluster moments. No bound involving the total ambient
volume is used. Mere finiteness of random all-order moment constants
would not suffice for the expectation arguments below.

For later use we recall why short stationary steps have the requisite
spatial sparsity. Conditional on a fixed transcript, each marginal is
reversible normally reflected Brownian motion, with variance parameter
two in every free coordinate. We use the symmetric Neumann
identification of the fixed-polytope process \cite{BassHsu} and the
form-domain forward--reverse identity \cite[Section 1, equation
(1.5)]{LyonsZheng}. The forward--reverse martingale
decomposition bounds each coordinate's displacement during $[0,T]$
by a fixed multiple of the suprema of its forward and reversed Brownian
martingales. Within either direction the coordinate martingales have
independent Brownian coordinates. If $k$ separated boxes have a
displacement exceeding $d>0$, select one of their finitely many shell
sites in each box; at least $k/2$ witnesses occur in one time direction.
The reflection principle and a union bound give
\begin{equation}\label{c:short-step}
 \Prob(\hbox{all these boxes move by more than }d)
 \le C\bigl(C_L e^{-c d^2/T}\bigr)^{k/2}.
\end{equation}
The constants are independent of the pinned values and of the number
of free coordinates. Fixed coordinates contribute no displacement.
Mix over transcripts. Combining \eqref{c:short-step} with initial edge
slack in \ref{f:conditioned-sparse}, and taking $T$ small after the slack parameter,
shows that transmission edges are confined to clusters of sparse boxes.
Connected-animal counting gives all cluster-size moments, their spatial
concentration, and arbitrarily small defect density by improving the
fixed parameters and shortening the testing step. Small constant drifts
are covered by the same reversible decomposition for their tilted
invariant measures, or by the bounded drift correction.

Here is the constant-drift extension of the finite-body identification.
After removing exact pins, let $K\subset\R^m$, $m\ge1$, be the
positive-volume polytope of free coordinates and put
\[
 \mu_b(\dd h)=Z_b^{-1}e^{b\cdot h}\,\dd h,
 \qquad \mathcal E_b(f,g)=\int_K\nabla f\cdot\nabla g\,\dd\mu_b.
\]
Its form domain is $H^1(K^\circ)$. An all-pinned fiber is deterministic
and requires no reflected process.
The density is smooth and bounded above and below on this fixed compact
body. Integration by parts for a smooth function $g$ with zero normal
derivative gives
\[
 \int_K f(\Delta g+b\cdot\nabla g)\,\dd\mu_b
       =-\mathcal E_b(f,g).
\]
Thus the weighted Neumann form identifies a stationary reversible
process with generator $\Delta+b\cdot\nabla$. For the coordinate $h_i$,
writing $\rho_b$ for the density and $n_{\rm out}$ for the outward unit
normal, the same calculation gives
\[
 \mathcal E_b(h_i,g)
 =-b_i\int_K g\,\dd\mu_b
  +\int_{\partial K}g\,n_{{\rm out},i}\rho_b\,\dd\mathcal H^{m-1}.
\]
The coordinate additive-functional argument in \cite{BassHsu} therefore
gives drift $b_i\dd t$ and an inward normal regulator. Its coordinate
martingale brackets are $2\delta_{ij}t$, since the energy density is
$2\nabla f\cdot\nabla g$. The coordinate functions belong to the form
domain, so the Lyons--Zheng identity applies even though they need not
satisfy the generator's Neumann boundary condition. Convex Skorokhod
uniqueness identifies this process with the reflected bath driven by
variance-two Brownian motion and drift $b$. This weighted-form argument
is separate from the zero-drift assertion of \cite{BassHsu}.

\subsection{Projection updates and square dissipation}

\begin{lemma}[Averaging structure and cluster displacement]\label{c:cluster}
On a zero-drift step, normal reflection of a difference acts on each
transmission cluster by compositions of averaging projections and their
convex mixtures. Free-input matrices are symmetric positive contractions
at each projection, with nonnegative entries and row and column sums
at most one. Pins supply their fixed difference; common individual
walls supply zero. For a tail with zero pinned input, insert coordinate
contractions after the averaging updates. If a cluster has $H$ free
vertices and total square loss $D$, then
\begin{equation}\label{c:displacement}
 \sup_{t\le T}\max_{x\in C}|u_t(x)-u_0(x)|
 \le C(H+1)\sqrt D.
\end{equation}
The same estimate holds for an untruncated conservative difference,
and is invariant under addition of a constant.
\end{lemma}
\begin{proof}
Use the finite-dimensional construction by projecting after each common
Brownian increment. Interpolate the two preprojection inputs, and also
their pinned values. Except at finitely many parameter breakpoints,
the derivative of projection is computed on a fixed active face.
An active edge makes its two free derivatives equal. A free component
therefore receives the average of its input derivatives, unless it
touches a pin or individual wall, when it receives the prescribed
derivative. Several pins active on a parameter interval have equal
derivatives. Integrating over the interpolation parameter gives the
asserted convex mixture of orthogonal averaging projections and zero
blocks. All intermediate contacts belong to the union of endpoint
near-contact clusters: an infinitesimal projection cannot contact an
edge whose slacks in both endpoint fields are bounded below.

For $u=(z-b)_+$, convexity and the pin bound $z|_P\le b$ give
$u'\le Au$ coordinatewise. Write this as $u'=RAu$ with a diagonal
matrix $0\le R\le I$. The lower tail is identical after changing sign.
Each elementary matrix $A$ or $R$ is a positive self-adjoint
contraction, hence
\[
 \|v-Av\|_2^2\le\|v\|_2^2-\|Av\|_2^2.
\]
Add zero as a possible initial level. Every intermediate coordinate
is a substochastic mixture of the initial levels, with lost mass put
at zero. Jensen's identity and the column deficits show that its
variance about that mixture is at most the total loss $D$. It is
therefore within $\sqrt D$ of some initial level. Each elementary
jump is at most $\sqrt D$. A gap greater than $4\sqrt D$ between
successive initial levels cannot be crossed. There are at most $H+1$
levels, which proves \eqref{c:displacement} uniformly in the projection
mesh. Passing to the reflected limit proves the assertion.
\end{proof}

We use $D_C=\sum_{x\in C}u_0(x)^2-\sum_{x\in C}u_T(x)^2$ in the
zero-drift case. For two tails, add their losses. With drift, $D_C$ denotes
the sum of the nonnegative losses at projection substeps; the separate
source terms are estimated at the end of the proof. General doubly
stochastic matrices cannot replace the projections in this lemma: a
permutation has no square loss.

\subsection{Accessible contacts in a fixed box}

The next construction supplies coercivity even when a nominated edge
inequality is redundant. Choose a fixed enlarged testing box, containing
at most $N$ free coordinates, and a collar $B$ separated from its central
test patch by the long-slack range. In the initial field use coordinates
$y_x=h_x-h_x(0)$. Adjoin a ground coordinate $y_*=0$ for all individual
constraints, and add the auxiliary bounds $|y_x|\le r_0$ on the collar.
These bounds restrict the paths we construct; they are not extra reflecting
walls of the original process.

\begin{lemma}[Facet paths and uniform forcing]\label{c:force}
Suppose every original edge involving a free coordinate has initial
slack at least $s>0$. Assume that in both directed orientations every
original-edge path between a central test vertex and the collar, without
an individual-constraint edge to ground, has residual slack greater
than $\Lambda>2$. Every free component is required to reach the collar
or an observed fixed vertex through original edges. Ordinary sign walls
are allowed, and a tested wall has initial slack at most one.

For every central edge there are at most $N+1$ genuine facet constraints
joining its endpoints, or grounding both endpoints at genuine pins or
walls. A tested individual wall is grounded by such a path. For any fixed
$T>0$ there are $p_0,I_0>0$, depending only on these fixed parameters,
such that every facet selected in these paths can receive normal impulse
at least $I_0$ in time $T$, on an event of the interior Brownian noises
of probability at least $p_0$, while all collar coordinates move by less
than $r_0$. Take $r_0<s/8$.
\end{lemma}
\begin{proof}
Write the inequalities in residual form $y_i-y_j\le c_{ij}$, with
individual inequalities using the ground vertex. All costs are
nonnegative because $y=0$ is feasible. Delete redundant inequalities.
The strongest bound on a difference in a finite system of difference
constraints is the shortest directed path cost in its constraint graph:
summing a path proves one direction, and shortest-path potentials give
the converse. The original central bound has cost at most two, so the
facet graph contains a path of cost at most two. Remove cycles. If it
uses ground, split the path there. Neither part can use an artificial
collar bound, since its segment from or to the tested endpoint would
already cost more than $\Lambda$. For a wall use the same argument
with budget one. A visit to a vertex carrying a wall is not itself a
visit to ground; only its individual-constraint edge gives ground input.

We justify uniformity of the forcing construction. Move every signed
coordinate a distance $\tau\le\min(s/8,r_0/4)$ toward its allowed side,
leaving other coordinates unchanged. All edge constraints retain slack
at least $s-2\tau$, and the artificial collar retains slack at least
$r_0-\tau$. A fixed-radius Euclidean ball about the new point lies in
the auxiliary polytope. Every feasible displacement satisfies
$|y_x|\le2N+r_0$, by following a path to a pin or collar. Individual
bounds with greater slack than $2N+r_0+1$ may be capped at that value:
they remain redundant. Thus, after fixing one of finitely many graph
and sign patterns, the effective parameters form a compact family
with a uniform interior ball.

Fix a parameter value and take the facet paths just obtained for its
polytope. For one facet choose a point in its relative interior where
all other distinct-normal inequalities are strict. The artificial
collar inequalities are among the strict ones. Drive the coordinates
from the initial point along a feasible segment to this point, and
then a fixed distance in the outward normal direction. The reflected
trajectory stays at the facet during the outward part and acquires
a positive impulse. An initial segment toward the interior-ball
center can be inserted. A sufficiently small uniform tube around
this piecewise linear driver preserves the collar inequality and at
least half of the final impulse.

For a fixed polytope, uniform-driver continuity and convergence of the
projection recursion are the deterministic Skorokhod assertions in
\cite[Section 2, Corollaries 3--4]{Slominski}. The geometric hypotheses
hold here: if $B(a,r)\subset K$ and $|a-x|\le R$ on $K$, every inward
unit supporting normal at $y$ has scalar product at least $r$ with
$a-y$; within distance $r/2$ of $x$ its scalar product with
$(a-x)/|a-x|$ is at least $r/(2R)$. Convexity supplies the exterior
sphere condition. On the final interval only one distinct facet
normal is active, and the one-dimensional running-minimum formula
gives impulse stability.

Parameter stability does not require interpreting this fixed-body
reference as a changing-body theorem. Fix the graph and the free
coordinates, and perturb initial heights and individual pin or wall
levels by at most $\delta$, in absolute height coordinates. Apply the
projection derivative argument of Lemma~\ref{c:cluster}, now allowing
these individual inputs to vary. Every output difference is a
subconvex combination of the previous free differences and the
individual-input differences. Under the same driver its supremum
norm is therefore at most $\delta$ at every projection step and in
the reflected limit. The original edge bounds remain exactly one.
Combining this parameter bound with fixed-body driver continuity
preserves the chosen tube in a parameter neighborhood. In the final
single-normal interval, normal balance also preserves the positive
impulse. Graph or free-coordinate patterns belong to a finite list
and are treated separately.

Brownian motion has positive probability of belonging to the chosen
tube: its piecewise linear center has square-integrable derivative,
so translation by that center preserves positivity of a Brownian
small-ball event. The construction persists in a neighborhood of
the parameter value. Merge identical free normals by retaining the
strongest bound. The selected relative-interior facet point then
moves slightly along its normal, with every other inequality still
strict. Use facets of each limiting polytope rather than following
a facet that disappears in a limit. A finite cover of the compact
parameter family supplies positive minimum probability and impulse.
\end{proof}

The long-slack assertion in \ref{f:sparse} verifies both orientations
required by Lemma~\ref{c:force}: residual costs telescope to path length
minus an oriented height difference. Absolute slack greater than two
excludes either direction of grounding through the far collar.
The uniform interior-ball argument deliberately excludes shrinking
two-sided intervals. Their initial edge-slack estimates can be uniform,
but this does not supply the full-dimensional forcing statement above.

\begin{lemma}[Insulation]\label{c:insulation}
Outside a sparse set of initially selected boxes, every interior forcing
event in Lemma~\ref{c:force} has conditional probability at least $p_0/2$
of occurring in the uncut process, given the initial pair and transcript.
\end{lemma}
\begin{proof}
Cut the edges across an initially strictly slack shell. For the cut
exterior process let $A_B$ be the event that an exterior endpoint moves
by more than $s/4$. If in the original stationary process every endpoint
of a cut edge moves by less than $s/4$, none of these edges is contacted.
Until a first cut contact the cut and uncut solutions agree. Uniqueness
therefore makes them agree throughout this quiet event. In particular
$A_B$ implies a shell-displacement event of the original process.
Equation~\eqref{c:short-step} gives joint sparse bounds for these events
on separated boxes. Lemma~\ref{c:conditional}, with $q=1/2$, shows that
boxes with $\Prob(A_B\mid\mathcal I)>1/2$ are still arbitrarily sparse.
On the remaining initial configurations exterior quietness has
probability at least $1/2$ and depends only on exterior noises. The
controlled interior noises are independent of them. Their intersection
has probability at least $p_0/2$, and no cut edge is contacted there.
Thus the cut controlled path is a path of the original process.
\end{proof}

\subsection{Coercivity and the cutoff estimate}

At an upper-edge contact of the first replica, $h_i-h_j=1$ and feasibility
of the second implies $z_i-z_j\ge0$. The contribution to the upper-tail
square loss is $2(u_i-u_j)\dd L\ge0$. The lower tail follows by reversing
the sign. Pin contacts and common walls have the same dissipative order.
Let $g$ be the initial jump at a facet selected in Lemma~\ref{c:force}.
It has a fixed bound in absolute value: use two for a lattice edge or
a zero-input pin edge, and the string height bound for an individual
wall. On the forced event, either a relevant coordinate moves by
at least $|g|/4$, or the jump retains magnitude at least $|g|/2$
during the impulse. Lemma~\ref{c:cluster} in the first case, and the
impulse in the second, give
\begin{equation}\label{c:force-loss}
 D_C\ge c g^2 H_C^{-m}
\end{equation}
for a fixed $m$. If the initial jump has the wrong sign for this facet,
the large-movement alternative must occur before contact. Forcing one
replica suffices; no force of the second replica is imposed.

To sum the local tests without counting a cluster loss once per vertex,
associate to a testing box $B$ the quantity
\[
 d_B=\sum_{C:\,C\cap B^+\ne\varnothing}\frac{D_C}{H_C},
\]
where $B^+$ is its fixed enlargement. A cluster meets at most $C_LH_C$
such boxes, so
\begin{equation}\label{c:allocation}
 \sum_Bd_B\le C_L\sum_CD_C.
\end{equation}
Let $H_B$ be the maximum incident cluster size in $B^+$. Increasing $m$
in \eqref{c:force-loss} absorbs the division by $H_C$. Conditional
Cauchy--Schwarz on the forced event $F$ gives
\[
 \E[d_B\mid\mathcal I]
 \ge \frac{c p_0^2 g^2}{1+\E[H_B^m\mid\mathcal I]},
 \qquad
 \Prob(F\mid\mathcal I)^2
 \le \E[H_B^m\mid\mathcal I]
       \E[\mathbf1_F H_B^{-m}\mid\mathcal I].
\]
Choose the edge and wall weights to be a fixed small multiple of
$[1+\E(H_B^m\mid\mathcal I)]^{-1}$. Telescoping along the finite facet
paths controls each tested original edge or anchor. Equation
\eqref{c:conditional-moments} supplies all required reciprocal moments,
and \eqref{c:allocation} supplies the expected coercivity inequality.

We now bound the dissipation without already assuming pointwise tightness
of $z$. Let $\eta$ be a deterministic cutoff, equal to one on an inner
box, zero outside an outer box. Let $d_{\rm lat}$ denote the gap
between those boxes in lattice units, so the cutoff has lattice
Lipschitz constant $C/d_{\rm lat}$.
For a cluster put $\eta_C=\max_C\eta$. Lemma~\ref{c:cluster} and
contraction give
\[
 \sum_{x\in C}|u_T(x)^2-u_0(x)^2|
 \le C H_C^m\sqrt{D_C}\,\|u_0\|_{\ell^2(C)}.
\]
Since $|\eta_x^2-\eta_C^2|\le
C\eta_C\operatorname{diam}(C)/d_{\rm lat}$, Young's inequality yields
\begin{equation}\label{c:cutoff-cluster}
 \sum_{x\in C}\eta_x^2(u_T(x)^2-u_0(x)^2)
 \le-\tfrac12\eta_C^2D_C
      +C d_{\rm lat}^{-2}H_C^m\sum_{x\in C}u_0(x)^2.
\end{equation}
Using the maximum cutoff avoids division by a possibly zero minimum
on clusters crossing the support boundary. Sum and take expectations;
stationarity cancels the left side. Cluster multiplicities and
diameters can be included in a larger fixed moment of $H_x$. Combining
with the force bound and using $d_{\rm lat}\asymp n(s-r)$ gives,
for $r<s<R$,
\begin{equation}\label{c:caccioppoli}
 \E\mathcal E_n(Q_r)+
 \E\sum_{C:\,C\cap Q_r\ne\varnothing}D_C
 \le C(s-r)^{-2}\E N_n(Q_s)+C,
 \qquad
 N_n(Q)=\int_{n,Q}H_x^m(u_0(x)^2+1).
\end{equation}
Testing enlargements are absorbed by an intermediate box.

For a free comparison choose at each time the constant minimizing
$\sum_x\eta_x^2(z_t(x)-c)^2$. The minimum at the next time is no larger
than the value at the previous minimizer. Its expectation is stationary,
so \eqref{c:cutoff-cluster} gives the same estimate. This step uses local
oscillations, bounded deterministically by $Cn$, rather than absolute
heights or the ambient diameter.

Here are the localization details relevant to that uniformity. On a
conservative cluster, constant cancellation and $|\nabla z|\le2$ give
\[
 |z_t(x)-z_0(x)|\le2\operatorname{diam}(C),\qquad
 D_C\le4H_C\operatorname{diam}(C)^2.
\]
In a cluster intersecting a cutoff support, compare every value with
one value in the support; the additional error is at most $2H_C$.
Thus the right side of \eqref{c:cutoff-cluster} is bounded by local
incident-size moments times local squared values plus a higher size
moment. Clusters reaching uncontrolled data across a margin of order
$n$ have negligible errors: their local contribution is bounded by
a constant times
\[
 \sum_{x\in Q_s}(nH_x+H_x^2)
       \mathbf1_{\{\operatorname{diam}(C_x)\ge c n\}},
\]
whose expectation tends to zero faster than any chosen negative power
by increasing the moment order. Alternatively, the marginal
forward--reverse decomposition gives ambient-uniform moments of
$\sup_{t\le T}|z_t(x)-z_0(x)|$; Holder with the crossing probability
gives the same result. Neither calculation introduces the number of
vertices outside the localization box.

\subsection{Absorbing the weighted square norm}

We have proved coercivity and \eqref{c:caccioppoli}. It remains to bound
the weighted norm on its right from macroscopic averages. On the
deterministic spatial-moment event, \ref{f:hole} supplies an extension
$U$ of a signed combination $F$ of the tails, with
\[
 \|\nabla U\|_{L^p}\le M\mathcal E_n^{1/2},\qquad
 \|W\|_{L^t}\le M,\qquad
 \|F-U\|_{L^{2t'}}\le M,
\]
where $W=H^m$, all norms are in scaled coordinates, and $M$ is
deterministic. Choose $t>p/[2(p-1)]$ and put
$\alpha=2-2/p-1/t>0$. Partition a fixed box into squares of side
$\ell$. Poincare--Sobolev on each square followed by Holder gives
\begin{equation}\label{c:interpolation}
 \int WF^2\le C\ell^{2\alpha}\mathcal E_n
 +C\ell^{-2/t}
   \left(1+\ell^2\sum_B|\operatorname{avg}_B F|^2\right).
\end{equation}
Indeed the oscillation estimate on one square is
$\|U-U_B\|_{2t'}\le C\ell^\alpha\|\nabla U\|_p$;
sum using $2t'>p$. For the piecewise constant averages use
$\sum_B\ell^2|U_B|^{2t'}\le
\ell^{-2(t'-1)}(\sum_B\ell^2|U_B|^2)^{t'}$.
The hole error is bounded by Holder with its stated moment.

For two tails take $F=(z-b)_+-(a-z)_+$. Their squares add to $F^2$,
and $F-z$ is bounded by a constant depending only on $a,b$.
Consequently \ref{f:averages} controls the averages in
\eqref{c:interpolation}. In the free version subtract the same local
average from $z$ and use equal-mass average differences. In a uniformly
bounded ordered comparison the weighted norm is already controlled.
Choosing $\ell$ after a desired $\eps>0$ yields
\begin{equation}\label{c:absorb}
 \E N_n(Q_r)\le\eps\E\mathcal E_n(Q_s)
                 +C\eps^{-m}(s-r)^{-k}
\end{equation}
with finite deterministic exponents. For instance $p=3/2,t=4$ gives
energy coefficient $\ell^{5/6}$ and average coefficient
$\ell^{-1/2}$. No higher energy moment is used. Exceptional events
are discarded using their arbitrarily fast polynomial probability,
local Lipschitz bounds, and the moments of one local average.

For completeness, combine \eqref{c:absorb} and
\eqref{c:caccioppoli} with an intermediate box. There exist finite
$A,m$ such that, for each $0<\theta<1$,
\[
 \E N_n(Q_r)\le\theta\E N_n(Q_s)
                  +C\theta^{-m}(s-r)^{-A}.
\]
Use successive gaps $d_j=d_02^{-j}$, with fixed
$\theta2^A<1$. The resulting series is bounded. Stop after
$J=\lfloor\kappa\log_2n\rfloor$ iterations, with $\kappa>0$ small
enough that every auxiliary interpolation box is still larger than
$n^{1/2}$. The terminal bound is polynomial in $n$, from local
Lipschitzness and one average moment. After fixing $\kappa$, decrease
$\theta$ so that $\theta^Jn^B\to0$ for that polynomial exponent $B$.
Choose the exceptional-event exponent last. This proves the local
bounded norm and \eqref{c:energy-bound}, and \ref{f:hole} then gives
the first assertion of \eqref{c:compactness} for every $1<p<2$.

\subsection{Rough one-sided observations}

A single lower tail need not have controlled averages. We instead use
the averages of $z$ to obtain a negative Sobolev bound. Here and below
the analytic norms and integrals are in scaled coordinates. Choose
a smooth cutoff $\chi_0$ supported in $\widehat Q_R$, equal to one
on a slightly smaller box containing the supports of the test
cutoffs used below. If $z^{\rm sc}(y)=z(ny)$ denotes the scaled
field with the fixed exterior convention of the experiment, write
\[
 \|z\|_{H^{-\eta}}
 :=\|\chi_0 z^{\rm sc}\|_{H^{-\eta}(\R^2)}.
\]
The dual space in the pairing is $H^\eta(\R^2)$; a test such as
$\chi^2U$ is extended by zero outside its compact support in the
region where $\chi_0=1$. Thus no boundary condition on $z$ at the
artificial box boundary is imposed by this norm. Dyadic smoothing,
\ref{f:averages}, and the bottom-scale Lipschitz bound give, for any
fixed $\eta>0$ on a compact box,
\begin{equation}\label{c:negative}
 \E\|z\|_{H^{-\eta}}^2
 \le C\sum_{j\ge0}2^{-2\eta j}(1+j)^2<\infty.
\end{equation}
The zeroth-scale average is included. Equal-mass difference bounds
alone would not suffice: the constants $z=-n$ have zero such
differences. The centered version uses its fixed local average in
this zeroth-scale term.

Let $u=(a-z)_+$ and let its extension $U$ copy a surrounding positive
tail value into each hole. If $D_x$ is a fixed multiple of the hole
diameter, the original Lipschitz bound gives pointwise
\begin{equation}\label{c:one-sided-pointwise}
 U^2\le(a-z)U+DU.
\end{equation}
For $U_x>0$, write $U_x=a-z_y$ at its copied rim site and use
$z_x-z_y\le D_x$; for $U_x=0$ the inequality is immediate.
Across a controlled observed boundary the original tail is zero,
and its zero extension is uniformly Lipschitz. Zero filling behind
such a boundary preserves \eqref{c:one-sided-pointwise}. An
uncontrolled seed is excluded by the stipulated positive margin.

Here is a quantitative absorption using only the first energy moment.
Choose $p=3/2$, $t=4$, and $\eta=1/6$. Let $\chi$ equal one on
$\widehat Q_r$, be supported in $\widehat Q_s$, and satisfy
$|\nabla\chi|\le C/(s-r)$. Put
$X=\|\chi U\|_2$, $Z=\|z\|_{H^{-1/6}}$, and $d=s-r$.
This $d$ is the dimensionless scaled gap; the corresponding lattice
gap is $nd$, as in the cutoff calculation above.
Pair \eqref{c:one-sided-pointwise} with $\chi^2$. Fractional
interpolation between $L^2$ and $W^{1,3/2}$ has exponent $1/4$
for $H^{1/6}$, while
\[
 \|\nabla(\chi^2U)\|_{3/2}
 \le C(\mathcal E_n^{1/2}+d^{-1}X).
\]
Duality and two applications of Young's inequality give
\[
 X^2\le\delta\mathcal E_n
 +C(1+\delta^{-1/4}+d^{-1/2})(1+Z^2).
\]
The derivative of $\chi^2U$ contains $\chi U$, which explains why
no uncontrolled larger-box $L^2$ term is introduced. Weighted Sobolev
interpolation to $L^{8/3}$, with exponent $3/8$, gives
\[
 \int_{\widehat Q_r}WU^2
 \le C\mathcal E_n^{3/8}X^{5/4}
       +C(1+d^{-3/4})X^2.
\]
Choose $\delta$ after $\eps$ and $d$ to obtain
\begin{equation}\label{c:one-sided-absorb}
 \int_{\widehat Q_r}WF^2\le\eps\mathcal E_n(Q_s)
            +C\eps^{-1}d^{-k}(1+Z^2),
\end{equation}
where $F=u$, and the filling error has been included by Holder.
One may take $k=5/4$ if the extension constants have no additional
gap dependence; only a finite polynomial exponent is used.
Take expectations using \eqref{c:negative}, and repeat the interior
iteration. This proves the one-sided part of the theorem. The
argument is annealed over the actual transcript by Lemma~\ref{c:freezing}.
To track the threshold, retain $a$ in
\eqref{c:one-sided-pointwise}: its pairing contributes at most
$|a|X$, and hence $C a^2$ after absorption. All later estimates are
linear in this lower-order bound. The weights depend on the original
fields and force tests, not on $a$. This proves
\eqref{c:threshold-scaling}. For $z/K$ the force impulse in its
difference is divided by $K$, while a facet jump is bounded by
$2/K$; thus the implication from linear impulse loss to the square
of that jump has the same constant. Division by $K\ge1$ also
improves the Lipschitz and negative-Sobolev bounds.

\subsection{Removing filled holes and obtaining microscopic tightness}

Fix first one hole construction. To prove its error tends to zero,
use a sequence of improved initial good parameters, indexed by $j$,
whose defective-edge densities are $\rho_j\downarrow0$. They may
use shorter forcing steps; the original hole construction is unchanged.
For every fixed $j$, the estimates already proved give good edges
$G_{j,n}$ and weights with
\[
 \E\sum_{e\in G_{j,n}}w_{j,e}|\nabla_eu|^2\le C_j,\qquad
 \E\int_n\mathbf1_{\{e\notin G_{j,n}\}}\le\rho_j+o(1),
\]
and low-weight density at most $C_j a^q+o(1)$ for $w_{j,e}<a$.
Since $|\nabla_eu|\le2$, split into defective edges, low-weight
edges, and the remaining edges. Cauchy--Schwarz on the last class gives
\begin{equation}\label{c:all-edge-l1}
 \E\int_n|\nabla_eu|
 \le2\rho_j+2C_j a^q+\frac{C_j^{1/2}}{n\sqrt a}+o(1).
\end{equation}
Send $n\to\infty$, then $a\downarrow0$, then $j\to\infty$.
Thus the annealed mean absolute jump over all edges tends to zero.

For holes of diameter at most a fixed $M$, telescope from each hole
site to its copied rim value. Paths and their edge congestion have
a fixed polynomial bound in $M$. The paths may use edges omitted
from the energy, since \eqref{c:all-edge-l1} now covers all edges.
For larger holes use the original Lipschitz bound and their uniformly
bounded high moments. Consequently
\[
 \E\int_{n,Q_r}|U_n-u|
 \le C M^k\E\int_{n,Q_s}|\nabla_eu|+C_q M^{-q}+o(1).
\]
First send $n\to\infty$ and then $M\to\infty$. This proves the
second part of \eqref{c:compactness}. The same argument applies to
the free centered difference. At a controlled observed boundary use
the Lipschitz zero extension of a tail. For common-wall comparisons
use the original difference on the full common graph; the wall
penalties are retained separately in \eqref{c:energy}.

The same decomposition yields the volume-tail bound
\[
 \Prob_{\rm vol}(n|\nabla_eu|>M)
 \le\rho_j+C_j a^q+\frac{C_j}{aM^2}+o(1).
\]
Choose $j$, then $a$, then $M$ to prove tightness. A finite path
gives each fixed-neighborhood difference. On a fixed step the
cluster displacement estimate and \eqref{c:energy-bound} give
\[
 \Prob_{\rm vol}\left(n\sup_{t\le T_0}|u_t-u_0|>M\right)
 \le\Prob_{\rm vol}(H>K)+\frac{C K^m}{M^2}+o(1).
\]
For an untruncated difference in a free region use the conservative
cluster estimate, or
$n\sup|z_t(x)-z_0(x)|\le
H_x\max_{e\subset C_x}n|\nabla_ez_0|$ and first restrict $H_x\le K$.
When $a<b$, the two tails do not control the part of $z$ inside
$[a,b]$; the assertion about $\nabla z$ therefore uses the free
centered comparison on a ball of radius a fixed positive multiple
of $n$ about the sampled region. A finite microscopic free
neighborhood is insufficient when nearby pinned differences
oscillate within $[a,b]$.
These are fixed-step tightness assertions, not an asserted modulus
of continuity for the $n$-rescaled process as $T\downarrow0$.

\begin{lemma}[Preserving a connected zero exterior]\label{c:exterior-zero}
Let $D_n$ be a finite simply connected triangular-lattice polygon.
Let $E_n$ consist of its boundary vertices and all vertices outside
its interior, and put $u=0$ on $E_n$. Assume $u$ is uniformly
Lipschitz on the full lattice and has the good-edge energy and
inverse-weight bounds established above, including edges incident
to its fixed boundary. The hole extension can be chosen so that
$U=0$ on $E_n$, with the same Sobolev estimate as in
\ref{f:hole}. Its normalized annealed $L^1$ error tends to zero
under the improved-parameter bounds used in
\eqref{c:all-edge-l1}. Thus the same-boundary comparison difference,
extended by zero, has a filled extension satisfying
\[
 \sup_n\E\|U_n\|_{W^{1,p}(\R^2)}^2<\infty,\qquad
 U_n=0\ \hbox{on }\R^2\setminus D_n
\]
after scaling and affine interpolation, for every $1<p<2$.
\end{lemma}
\begin{proof}
Use the original lattice vertices for the extension. Bad boxes are
first enlarged by the fixed amount needed to control every original
edge on a surrounding good rim. Their outer filled hulls give the
holes of \ref{f:hole}; passing from boxes to their unions of lattice
sites changes diameter and counting constants only by the fixed box
size. Enlarge the rim by a fixed layer if needed to include an exit
by a triangular-lattice diagonal edge. The same square-coordinate
contour construction supplies connected good rim paths of length
at most $C(d+1)^2$ for a hole of diameter $d$. No extra probabilistic
claim for independently marked individual sites is needed.

The zero set $E_n$ is connected to infinity in the lattice graph.
Indeed the polygon has one connected boundary cycle, and its
exterior is the unbounded complementary component; including
boundary vertices removes any vertex-level corner separation.
For a hole $H$ which meets $E_n$, follow an $E_n$-path from a
vertex in $H$ to infinity. Since $H$ is finite, the path exits it.
Its first exterior rim vertex lies in $E_n$ and has value zero.
Choose this vertex as the anchor of $H$, and fill all of $H$
with zero. If $H$ does not meet $E_n$, choose any surrounding good
anchor as before. Outside the holes retain the original values.
The resulting field is zero on every vertex of $E_n$.

We verify the energy bound for this choice, rather than relying
on preservation by an arbitrary anchor. If an edge joins a good
vertex $x$ to a hole, its new jump is the difference of $u(x)$
and the value at that hole's chosen rim anchor. Join these two
rim vertices by the good rim path just described and telescope.
The path length, the number of incident edges, and the repeated
uses of a rim edge have the same polynomial diameter bounds for
every choice of anchor. Applying the path inequality and then
Holder gives exactly the weighted bound in \ref{f:hole}, with
possibly a changed fixed exponent of $1+d$. Edges with both
endpoints in a filled hole have zero jump. Edges outside the
holes retain their old jump; exterior zero--zero edges may be
given weight one and contribute no energy. In particular the
weighted gradient estimate holds across the actual polygon
boundary. The spatial moment estimates and expected energy
bound then give its squared $W^{1,p}$ gradient bound. The local
weighted square bound already proved controls its $L^p$ norm.

The new anchor also preserves the vanishing filling error.
For holes of diameter at most $M$, join each vertex to its
selected anchor by a full-lattice path of length at most $CM$.
Such paths lie in a fixed enlargement of the hole's bounding
box, and their congestion is bounded by a fixed polynomial of
$M$. The all-edge estimate \eqref{c:all-edge-l1} applies to
the zero-extended original field, including boundary-incident
edges. It makes the expected error from these holes tend to
zero for each fixed $M$. For larger holes,
$|U(x)-u(x)|\le C(1+d_x)$ by full-lattice Lipschitzness; high
diameter moments make their error tend to zero as $M\to\infty$.
These estimates also cover the exceptional event by setting
$U=0$ there and using the established moment bounds.

For a comparison with identical boundary heights,
$z=h-\widetilde h=0$ on the boundary. Its extension by zero
is $2$-Lipschitz on the full lattice: edges crossing the
boundary first meet a fixed boundary vertex, and exterior
edges have zero difference. Apply the construction to $z$,
or to both of its tails and subtract. All triangle vertices
outside the lattice polygon have value zero, so affine
interpolation vanishes there as well. Norm equivalence on
the fixed lattice triangles gives the displayed global
Sobolev bound. This proof uses the unbounded connected
exterior, and does not apply to a hole containing an entire
isolated internal component of observed vertices.
\end{proof}

\subsection{Small drifts and completion of the proof}

In the projection scheme a relative drift of supremum norm
$c n^{-2}$ is inserted between averaging updates. The substochastic
matrices do not increase its supremum norm, so its transported
contribution during one step is at most $cTn^{-2}$. The displacement
bound acquires only a polynomial cluster factor times this quantity.
At the square-energy identity the additional source term is bounded by
\[
 C n^{-2}\int_0^T\sum_x\eta_x^2(|u_t(x)|+1)\dd t
 \le C n^{-2}\int_0^T\sum_x\eta_x^2(u_t(x)^2+1)\dd t.
\]
After stationarity this has the same form as the cutoff error, since
the lattice gap is at most $Cn$. The force test acquires an error
$C H^m n^{-4}$ in its square bound. Its spatial sum is bounded, and
indeed tends to zero on the normalized scales in use. The preceding
coercivity, absorption, and localization therefore remain valid.

It remains to check the change of marginal law for a smooth tilt.
Its density is proportional to
$\exp(n^{-2}\sum_x b(x/n)h(x))$, with bounded support at scale $n$.
For sufficiently small $\|b\|_\infty$, the exponential average
bound in \ref{f:averages} bounds a fixed moment greater than one
of this density and its normalizing factor. Holder transfers any
sparse event probability $C\rho^k$ to $C'\rho^{c k}$, and an
arbitrarily fast polynomial exceptional bound to another such
bound. It also transfers the finitely many average and cluster
moments used above. This argument is applied before adaptive
revealing; Lemma~\ref{c:freezing} then preserves the tilted marginals.
There is no claim of a uniformly bounded density for each exact
transcript. These observations complete the proof of
Theorem~\ref{c:comparison}.

\begin{remark}\label{c:scope}
The dependencies \ref{f:conditioned-sparse}, \ref{f:hole}, and \ref{f:averages}
must hold with their stated bounded-data and sign-string scope.
The zero-pinned controlled-gradient theorem alone does not give
that scope. In particular initial edge strictness is established
there by a conditional cone filling and homothetic volume estimate,
before Lemma~\ref{c:force} is used. The present proof also uses the
deterministic spatial moment bounds outside arbitrarily rare events,
not merely annealed finiteness of moment constants. No assertion
about uniform forcing on shrinking two-sided wall polytopes is made.
\end{remark}
